\documentclass[12pt]{report}
\usepackage[utf8]{inputenc}
\usepackage[czech]{babel}
\usepackage{hyperref}
\usepackage{graphicx} 
\usepackage{float}
\usepackage{enumitem}


\title{Překladač vlastního jazyka do instrukcí PL/0}
\author{
  Kotva, Pavel\\
  \texttt{kotva@students.zcu.cz}
  \and
  Kohlíček, Jan\\
  \texttt{kohl@students.zcu.cz}
}

\begin{document}

\begin{titlepage}
\begin{flushleft} 
{\includegraphics[width=.5\textwidth]{./img/fav_logo.jpg}\\[3cm]}
\end{flushleft}
\begin{center}

{\Large KIV/FJP - Semestrální práce}\\
\vspace{0.3cm}
{\Huge Překladač vlastního jazyka\\do instrukcí PL/0}\\
\vspace{1cm}
{\large Pavel Kotva - A17N0077P}\\
\vspace{0.1cm}
{\normalsize kotva@students.zcu.cz}\\
\vspace{0.7cm}
{\large Jan Kohlíček - A17N0075P}\\
\vspace{0.1cm}
{\normalsize kohl@students.zcu.cz}
\vfill
{\large \today}
\end{center}
\end{titlepage}

\tableofcontents
\thispagestyle{empty}

\chapter{Zadání}
Cílem práce bude vytvoření překladače zvoleného jazyka. Je možné inspirovat se jazykem PL/0, vybrat si podmnožinu nějakého existujícího jazyka nebo si navrhnout jazyk zcela vlastní. Jazyk bude překládán do instrukcí PL/0. 
\\
\\
Jazyk musí mít základní konstrukce:
\begin{itemize}[noitemsep]
	\item [-] definice celočíselných proměnných
	\item [-] definice celočíselných konstant
	\item [-] přiřazení
	\item [-] základní aritmetiku a logiku (+, -, *, /, AND, OR, negace a závorky, operátory pro porovnání čísel)
	\item [-] cyklus (libovolný)
	\item [-] jednoduchou podmínku (if bez else)
	\item [-] definice podprogramu (procedura, funkce, metoda) a jeho volání
\end{itemize}
\ \\
Rozšiřující konstrukce:
\begin{itemize}[noitemsep]
	\item [-] každý další typ cyklu
	\item [-] datový typ boolean a logické operace s ním
	\item [-] else větev
	\item [-] násobné přiřazení
	\item [-] paralelní přiřazení
	\item [-] parametry předávané hodnotou
	\item [-] návratová hodnota podprogramu
\end{itemize}

\chapter{Syntaxe}

\section{Proměnné a přiřazení}
Jazyk je staticky typový, všechny proměnné musíme nejprve deklarovat s jejich datovým typem.\\
\\
Datové typy:
\begin{description}
	\item [int] - označení celočíselného datového typu.
	\item [boolean] - logický datový typ. Boolean nabývá dvou hodnot: \textbf{true} (pravda) a \textbf{false} (nepravda). 
\end{description}
\ \\
Deklarace proměnné bez přiřazení:
\begin{verbatim}
int x;
boolean y;
\end{verbatim}
\ \\
Deklarace proměnné s přiřazení:
\begin{verbatim}
boolean x := true;
boolean y := false;
int z := (3 * 3) + 1;
\end{verbatim}
\ \\
Deklarace proměnných s násobným přiřazením:
\begin{verbatim}
int a;
int b;
int c;
a, b, c := 5;
\end{verbatim}
\ \\
Deklarace proměnných s paralelním přiřazením:
\begin{verbatim}
int a;
int b;
[a, b] := [10, 15];
\end{verbatim}





\subsection{Konstanty}
Konstanty jsou deklarovány s klíčovým slovem \textbf{const} a mohou obsahovat jen celé číslo.
\begin{verbatim}
const KONSTANTA := 5;
\end{verbatim}






\section{Podmínky (větvení)}
Podmínky zapisujeme pomocí klíčového slova \textbf{if}, za kterým následuje logický výraz. Pokud je výraz pravdivý, provede se následující příkaz. Pokud ne, následující příkaz se přeskočí a pokračuje se větví \textbf{else}. Zapsat větev \textbf{else} je povinné.
\begin{verbatim}
int i := 1;

if(i > 1)
start
    i := 2;
end
else
start
    i := 4;
end
\end{verbatim}


\subsection{Operátory}
Relační operátory, které můžeme ve výrazech používat:

\begin{table}[H]
\centering
\begin{tabular}{ c |  p{6cm}}
 Operátor & Zápis \\ \hline\hline
 	
Rovnost & \textbf{=} \\ \hline  
Je ostře větší & \textbf{\textgreater} \\ \hline
Je ostře menší & \textbf{\textless} \\ \hline
Je větší nebo rovno & \textbf{\textgreater =} \\ \hline
Je menší nebo rovno & \textbf{\textless =} \\ \hline
Nerovnost & \textbf{!=} \\ \hline
Obecná negace & \textbf{!} \\ \hline
A zároveň & \textbf{\&\&} \\ \hline
Nebo & \textbf{\textbar \textbar}  \\ \hline
\end{tabular}
\caption{Relační operátory}
\label{table:relacni_operatory}
\end{table}
\ \\
Podmínky je možné skládat a to pomocí dvou základních logických operátorů:
\begin{table}[H]
\centering
\begin{tabular}{ c |  p{6cm}}
 Operátor & Zápis \\ \hline\hline
A zároveň & \textbf{\&\&} \\ \hline
Nebo & \textbf{\textbar \textbar}  \\ \hline
\end{tabular}
\caption{Logické operátory}
\label{table:logicke_operatory}
\end{table}

\section{Cykly}


\subsection{While}
Prvním typem cyklu je \textbf{while} cyklus funguje jednoduše opakuje příkazy v bloku dokud platí podmínka. Syntaxe cyklu je následující:

\begin{verbatim}
int i := 0;

while(i < 4)
start
    i := i + 1;
end
\end{verbatim}

\subsection{Do\textendash while}
Posledním typem cyklu je \textbf{do\textendash while}. Je téměř stejný jako while, ale kontrolní podmínka je umístěna až na konec cyklu. Máme tedy jistotu, že minimálně jednou cyklus vždy proběhne. Syntaxe cyklu je následující:

\begin{verbatim}
int i := 0;

do
start
    i := i + 1; 
end while(i < 4)
\end{verbatim}


\section{Procedury}
Procedura je logický blok kódu, který jednou napíšeme a poté ho můžeme libovolně volat bez toho, abychom ho psali znovu a opakovali se. Proceduru deklarujeme v globálním prostoru, někde nad \textbf{main}.

\subsection{Procedura bez parametrů}
Syntaxe procedura bez parametrů je následující:
\begin{verbatim}
procedure nazevprocedury()
start
    int b := 5;
end

call nazevprocedury();
\end{verbatim}


\subsection{Procedura s parametry}
Procedura může mít také libovolný počet vstupních parametrů (někdy se jim říká argumenty), které píšeme do závorky v její definici. Rozšíříme tedy stávající proceduru o parametr \textbf{int a} a ten potom přidáme s konkrétní hodnotou do volání procedury:
\begin{verbatim}
procedure nazevprocedury(int a)
start
    int b := 5 + a;
end

call nazevprocedury(3);
\end{verbatim}

\subsection{Procedura s návratovou hodnotou}
Procedura může dále vracet nějakou hodnotu, správně by se ji mělo říkat funkce, ale z historických důvodů a zachování kompatibility nazýváme procedura.  Procedura může vracet právě jednu hodnotu pomocí příkazu \textbf{return}. Syntaxe je následující:
\begin{verbatim}
procedure nazevprocedury(int a)
start
    int b := 5 + a;
end
return b;

int y;
call nazevprocedury(3)(y);
\end{verbatim}




\chapter{Implementace}
Pro implementaci překladače byl vybrán jazyk \textbf{Java} a nástroj pro tvorbu syntaktických analyzátorů, kompilátorů a překladačů gramatiky \textbf{ANTLR}. \textbf{ANTLR 4.7} byl vybrán pro jeho snadnou a rychlou použitelnost.
\\\\
Popis adresářové struktury:
\begin{itemize}[noitemsep]
\item [+] \texttt{docs} - dokumentace
\item [+] \texttt{src} - zdrojové kódy
\item [.] \texttt{FJPLexer.g4} - obsahuje tokeny
\item [.] \texttt{FJPParser.g4} - obrahuje gramatiku
\end{itemize}


\chapter{Vzorové programy a generované instrukce}

\section{Program: Podmínky}
Zdrojový kód:
\begin{verbatim}
kod
\end{verbatim}
\ \\
Instrukce:
\begin{verbatim}
instrukce
\end{verbatim}


\section{Program: Cykly}
Zdrojový kód:
\begin{verbatim}
kod
\end{verbatim}
\ \\
Instrukce:
\begin{verbatim}
instrukce
\end{verbatim}


\section{Program: Procedury}
Zdrojový kód:
\begin{verbatim}
kod

\end{verbatim}
\ \\
Instrukce:
\begin{verbatim}
instrukce

\end{verbatim}


\chapter{Závěr}
V semestrální práci byl vytvořen překladač do instrukcí \textbf{PL/0}. Splnily jsme všechny body, které jsme si určily v zadání.
Syntaxe je originálním mixem syntaxe \textbf{C} a \textbf{Pascalu}.
Zdrojové kódy jsou k dispozici na githubu: \url{https://github.com/Ahantuon/fjp} 


\end{document}